%% ============================================================
%% Chapter 17: Sparse Recursive Holographic Steganography
%% ============================================================

\chapter{Sparse Recursive Holographic Steganography}
\label{ch:ch17-sparse-holographic-steganography}
\section{Introduction}

Chapters~\ref{ch:ch15-compression-fallacy} and~\ref{ch:ch16-compression-after-expansion} examined representations whose brevity conceals the machinery, or the exploration history, that makes them work. This chapter examines a representation problem of a different shape: not how to compress a payload, but how to distribute it so that no single location in a carrier is uniquely responsible for any fragment of it, and partial observation degrades gracefully rather than failing at a fixed boundary. The architecture developed here, Sparse Recursive Holographic Steganography (SRHS), replaces the standard one-pass map from message bits to carrier coefficients with three coupled operations: holographic payload projection, sparse carrier routing, and recursive residual correction.

Steganographic systems are conventionally organized around a spatial metaphor. A payload is partitioned into symbols; symbols are assigned to pixels, blocks, tokens, frames, or transform coefficients; a decoder later reads those assignments back. This organization produces three linked weaknesses. Locality produces brittle failure, since destroying the region assigned to a payload block destroys that block regardless of how much unused evidence remains elsewhere in the carrier. A one-pass encoder commits to its full modification before it can observe the realized consequences of any of it, and carrier channels are heterogeneous enough that post-embedding transformations are only partially predictable in advance. Dense modification, finally, wastes distortion on locations that contribute little reliable information while enlarging the statistical surface available to a detector.

The architecture proposed here replaces symbol placement with field construction. The payload is projected into a set of correlated measurements, each individually insufficient but jointly informative about the payload as a whole. These measurements are embedded into a sparse subset of carrier regions. After each embedding step the system decodes its own intermediate output, estimates the residual payload error, and selectively revises the carrier. Computation is therefore recursive, while modification remains sparse at every step. The term \emph{holographic} is used operationally rather than by analogy to physical optics: it denotes distributed recoverability, in which partial observations preserve degraded evidence about the whole message rather than exact evidence about isolated fragments of it. The term \emph{recursive} likewise denotes a concrete recurrence with shared parameters and an explicit stopping rule, not literal self-similarity or unbounded depth.

\section{Problem Formulation}

Let a cover object be $x \in \mathcal{X}$, drawn from a declared cover distribution $\mathcal{D}_\mathcal{X}$; let a payload be $m \in \{0,1\}^L$; and let a secret key be $k$, drawn from a keyspace large enough to resist exhaustive search. An encoder produces $y = E_\theta(x,m,k;B)$, where $y$ is the stego object and $B$ is a hard distortion budget. A channel transformation $T \sim \mathcal{T}$ may crop, compress, rescale, transcode, paraphrase, drop packets, reorder frames, or otherwise alter the object before it reaches the receiver. The decoder returns $\hat m = D_\phi(T(y),k)$. A useful steganographic system must jointly minimize decoding error, perceptual or semantic distortion, and detectability. For a detector $A$, define its advantage under equal priors as $\operatorname{Adv}(A) = |\Pr[A(y)=1] - \Pr[A(x)=1]|$. Every result below is stated relative to a declared $(\mathcal{D}_\mathcal{X}, \mathcal{T}, B)$ triple, and no claim is intended to hold outside such a declaration, a discipline this chapter shares with the falsifiable-claims requirement Chapter~\ref{ch:ch14-model-capsules} imposed on the deployment-native training architecture.

\section{Holographic Payload Projection}

A localized representation maps payload block $m_i$ primarily to a single carrier region $j$. SRHS instead maps the entire payload to a redundant collection of keyed measurements. An outer error-correcting code first produces $c = C(m) \in \{0,1\}^n$, centered as $\bar c = 2c - \mathbf{1} \in \{-1,+1\}^n$ to remove a codeword-dependent DC component, and a keyed linear projection generates
\[
z = \frac{1}{\sqrt n}\, P_k \bar c, \qquad P_k \in \mathbb{R}^{q \times n},
\]
followed by a bounded quantization $u = Q(z)$. The $1/\sqrt n$ normalization keeps the expected measurement energy independent of $n$. The rows of $P_k$ are generated from a cryptographic pseudorandom seed and normalized so evidence is dispersed across measurements rather than concentrated in a few of them.

\begin{proposition}[Restricted preservation]
\label{prop:rip}
There exists a family of subsets $S \subseteq \{1,\ldots,q\}$, occurring with non-negligible probability under the transformations in $\mathcal{T}$, such that the restricted projection $P_{k,S}$ preserves the pairwise distances among codewords of $C$ up to a bounded distortion factor $\delta(S) \in [0,1)$: for any two codewords $c_a \neq c_b \in C$ with centered forms $\bar c_a, \bar c_b$,
\[
(1-\delta(S))\,\|\bar c_a - \bar c_b\|^2 \leq \|P_{k,S}(\bar c_a - \bar c_b)\|^2 \leq (1+\delta(S))\,\|\bar c_a - \bar c_b\|^2.
\]
\end{proposition}

This is a restricted-isometry-type condition, applied here to a finite codeword set rather than to arbitrary sparse vectors. Since $\bar c = 2c - \mathbf{1}$, codewords at Hamming distance $d$ satisfy $\|\bar c_a - \bar c_b\|^2 = 4d$, so the minimum squared distance among centered codewords is $4d_{\min}$ for the outer code's minimum Hamming distance $d_{\min}$. Under Proposition~\ref{prop:rip}, the minimum squared distance among the normalized measurement images is at least $(1-\delta(S)) \cdot 4d_{\min}/n$, and nearest-codeword decoding over the surviving measurements $S$ succeeds whenever the aggregate noise satisfies $\|\eta_S\| < \sqrt{(1-\delta(S))\, d_{\min}/n}$. This margin gives a concrete design target: $\delta(S)$ can be estimated empirically on surviving subsets under the declared transformation distribution and compared against $d_{\min}/n$, a falsifiable numerical claim rather than an appeal to the assumption's plausibility.

Under this construction, the condition for recovery is not that every crop reconstruct the message exactly. It is that a sufficiently large family of measurement subsets retains enough rank, or a restricted-isometry-type property, to support decoding. For multimedia carriers, a view is any transformation that preserves a subset or mixture of carrier evidence: a crop, a temporal window, a frequency band, a resolution level, or a set of surviving packets; training samples multiple views of the same stego object and requires each to predict the same payload, the learned counterpart of the projection's algebraic redundancy.

\section{Sparse Recursive Embedding}

Divide the carrier representation into $N$ addressable regions, which may be spatial patches, wavelet bands, video tubes, audio time-frequency cells, text spans, or learned latent groups. At recursion step $t$ the system maintains a stego candidate $y_t$, a residual payload state $r_t$, and a remaining budget $b_t$, initialized as $y_0 = x$, $r_0 = u$, $b_0 = B$. A controller with shared parameters $\psi$ computes region scores $s_t = G_\psi(F(y_t), r_t, b_t, k)$, where $F$ is a carrier feature extractor. A sparse router selects an active set $S_t$ of at most $K \ll N$ regions; a proposal network produces a bounded revision $\Delta_t = U_\theta(F(y_t), r_t, k, S_t)$, and the carrier is updated only on the active set:
\[
y_{t+1} = \Pi_{\mathcal{C}(x,B)}\bigl(y_t + M(S_t) \odot \Delta_t\bigr),
\]
where $M(S_t)$ is the active-region mask and $\Pi_{\mathcal{C}(x,B)}$ projects the candidate back into the admissible distortion set around $x$. The system then performs an internal decode, $\tilde u_{t+1} = D_\phi^{\mathrm{inner}}(y_{t+1}, k)$, $r_{t+1} = u - \tilde u_{t+1}$. Recursion terminates when the estimated robust decoding loss falls below a threshold, when the budget is exhausted, or when a fixed maximum depth $R$ is reached. Because $G_\psi$, $U_\theta$, and the internal decoder share parameters across steps, increasing depth increases adaptive computation without increasing parameter count. Each step is simultaneously an intervention and a measurement: the encoder does not merely add signal to the carrier, it learns from what survived its own preceding attempt.

Sparsity without recursion commits irrevocably to a single, possibly mistaken, active set. Recursion without sparsity perturbs the entire carrier at every step and compounds detectable residue across steps. Their conjunction permits exploration followed by correction: early steps allocate weak measurements to regions predicted to be safe, an intermediate decode reveals which payload directions remain uncertain, and later steps address those directions through new regions or revisions of unreliable ones.

\begin{definition}[Marginal routing value]
Let $\ell_t$ denote expected payload loss after step $t$ and $c_{t,j}$ the cost of revising region $j$ at step $t$. The marginal value of revising region $j$ is
\[
v_{t,j} = \frac{\mathbb{E}[\ell_t - \ell_{t+1} \mid j]}{c_{t,j} + \lambda\, d_{t,j}},
\]
where $d_{t,j}$ estimates the increment to detectability contributed by revising $j$.
\end{definition}

Sparsity is a consequence of most regions carrying low or negative marginal value at a given state, rather than a constraint imposed for its own sake. Sparse recursion additionally supports conditional depth: covers with low channel noise can stop early, while difficult carriers consume additional steps, converting embedding computation from a fixed per-object expense into an allocation problem solved per example.

\section{Training Objective}

Training jointly optimizes the projection-aware decoder, the sparse recurrent encoder, and a population of steganalyzers. For sampled transformations $T_1,\ldots,T_J \sim \mathcal{T}$, the payload loss is $\mathcal{L}_{\mathrm{msg}} = \frac{1}{J}\sum_j \operatorname{BCE}(m, D_\phi(T_j(y_R),k))$. The fidelity term combines a domain metric with a feature-space discrepancy, $\mathcal{L}_{\mathrm{fid}} = d_\mathcal{X}(x,y_R) + \alpha\|\Phi(x)-\Phi(y_R)\|_2^2$. Security uses an ensemble of detectors, each trained to maximize discrimination, $\mathcal{L}_{A_i} = -\mathbb{E}_x \log(1-A_i(x)) - \mathbb{E}_{y_R}\log A_i(y_R)$, while the encoder is trained against the current ensemble using the non-saturating confusion loss $\mathcal{L}_{\mathrm{sec}} = -\frac{1}{M}\sum_i \log(1-A_i(y_R))$, which avoids the vanishing-gradient regime the saturating form produces early in training. Sparse, non-repetitive routing is encouraged by an $\ell_0$-style penalty on the gate vector plus a cross-step redundancy penalty, and a view-consistency term requires partial views to agree on the decoded payload distribution via a symmetrized KL divergence. Training alternates two minimizations, one over the encoder and decoder parameters jointly with the security and consistency terms, and one over the detector ensemble, rather than a single min-max expression. The hard distortion constraint is enforced by the projection $\Pi_{\mathcal{C}(x,B)}$ regardless of whether the learned penalties are well calibrated, which prevents the optimizer from trading visible carrier damage for payload accuracy under an imperfectly weighted objective.

\section{Analysis}

\subsection{Graceful Erasure Recovery}

Because $S$ is a discrete, bounded index set, recovery probability as a function of $|S|$ cannot be continuous in the analytic sense, and bounded-noise decoding can exhibit genuine thresholds at the outer code's correction radius. The claim is stated instead as monotone expected degradation, which is what the construction actually delivers.

\begin{proposition}[Monotone degradation under erasure]
\label{prop:erasure}
Suppose Proposition~\ref{prop:rip} holds and the projected codeword measurements are embedded with independent bounded noise. Let $S$ be a random subset of surviving measurements under a sampling model in which each measurement survives independently with probability $\pi$, and let $\mathrm{err}(S)$ denote the expected outer-decoding error conditional on $S$. Then $\mathbb{E}_S[\mathrm{err}(S)]$ is monotonically non-increasing in $\pi$, and no fixed payload block is dropped as a deterministic function of a predetermined carrier boundary, in contrast to a block-local scheme in which each payload block fails outright whenever its assigned region is erased, regardless of how much unused evidence survives elsewhere in the carrier.
\end{proposition}

Restricting the projection to a larger surviving set weakly improves the preserved codeword distances, since additional independent measurements can only add information to the linear system, so $\delta(S)$ is non-increasing in $|S|$ in expectation under the sampling model, and the recovery threshold is non-decreasing as $|S|$ grows. Increasing $\pi$ stochastically dominates the resulting $|S|$, so the probability of meeting the threshold, and hence the expected error, inherits monotonicity in $\pi$ by composition. This result is conditional rather than absolute: an adversary who removes all informative regions, destroys the carrier outright, or estimates the key-dependent projection structure can still prevent recovery. Holographic distribution changes the shape of the failure curve; it does not remove the underlying channel-capacity limit.

\subsection{Recursive Dominance Under Uncertain Response}

Weak dominance of a two-step encoder over a one-step encoder follows immediately from policy-class containment: any feasible one-step policy is available to the two-step encoder as the choice of a zero second update, so the two-step optimum cannot be worse. Strictness is a value-of-information claim. Model the realized carrier response as an unobserved state $h$, drawn from a prior unknown to the encoder before any modification is rendered, and let the first-step update produce an observation $o_1$ informative about $h$.

\begin{proposition}[Strict dominance via value of information]
\label{prop:recursion}
The two-step encoder's optimal expected loss is $\mathbb{E}_{o_1}[\min_a \mathbb{E}[\ell(a,h) \mid o_1]]$, and the one-step encoder's optimal expected loss is $\min_a \mathbb{E}[\ell(a,h)]$. The former is no greater than the latter in general, and strictly smaller whenever $o_1$ has positive value of information for the allocation decision.
\end{proposition}

The inequality holds in general because the right-hand side is the loss of a single action chosen without conditioning on $o_1$, feasible but not necessarily optimal for each realization of $o_1$; taking the expectation of the pointwise minimum cannot exceed the minimum of the expectation. Equality holds exactly when the optimal action under the prior remains optimal under every realization of $o_1$, that is, when $o_1$ carries no information relevant to selecting the action. This formalization makes explicit what recursion cannot buy: if the internal decode is uninformative about the realized channel response, for instance because the channel is deterministic and already known to the encoder, the value of information is zero and additional recursion steps yield no improvement in expectation, regardless of added computation. This is the discrete-communications analogue of Chapter~\ref{ch:ch10-clio-architecture}'s null-intervention principle: recursion, like suspension, is warranted only when the intervening observation is actually informative, and both frameworks refuse to treat additional computation or additional caution as valuable by default.

\subsection{Sparse Modification and Detectability}

Sparsity does not automatically imply security. A large change concentrated in a few regions can be easier to detect than a small change spread across many. The governing quantity is the divergence between the cover and stego distributions under the realized modification policy, not the cardinality of the active set as such. Sparse routing improves the frontier only when the controller correctly identifies high-capacity regions and respects per-region amplitude constraints.

\begin{proposition}[Capacity bound]
\label{prop:capacity}
Let $C_\mathcal{T}(B)$ denote the effective capacity of the transformed carrier channel under distortion budget $B$. Reliable transmission of an $L$-bit payload requires, asymptotically and subject to the standard qualifications of channel coding, $L \leq C_\mathcal{T}(B)$.
\end{proposition}

Projection redundancy, recursion, and learned routing cannot violate this bound; their function is to move the achievable operating point closer to it by using the available capacity more effectively, not to relax the bound itself.

\section{Experiments and Falsifiable Claims}

The principal comparison holds cover distribution, payload length, key length, distortion budget, training transformations, and compute reporting constant across conditions, comparing SRHS against traditional transform-domain embedding, a strong one-pass learned encoder, a dense recurrent encoder, and a sparse non-recurrent encoder. Primary outcomes are bit error rate after transformation, exact-message recovery rate, detector AUC, calibrated distortion, and encoding cost, together with the empirical distortion factor $\delta(S)$, estimated directly and compared against the decoding margin threshold, testing whether Proposition~\ref{prop:rip} actually holds at the operating point used elsewhere in the evaluation rather than being assumed to hold because the rest of the pipeline appears to work.

The decisive ablations remove one component at a time: holographic projection, outer error correction, sparse routing, recurrent residual correction, view consistency, and adversarial security training, with a further control replacing adaptive routing with randomly selected regions matched for sparsity and distortion. Three preregistered hypotheses define success: SRHS should lower transformed-channel bit error without increasing held-out detector AUC at matched distortion and payload; partial-view recovery should degrade smoothly with retained evidence, consistent with Proposition~\ref{prop:erasure}; and recursion should improve the frontier beyond a parameter-matched one-pass model, consistent with Proposition~\ref{prop:recursion}. Failure is equally specified in advance: if gains disappear against detectors not used in training, the method has learned the training steganalyzers rather than steganographic security; if the projection helps only because it adds redundancy, a conventional code at matched rate is sufficient; if additional recursion steps offer no gain after compute matching, adaptive revision is unnecessary; if robustness is obtained only by exceeding the declared distortion budget, the central claim of the paper is false.

\section{Security Model and Misuse}

Three observers are distinguished, extending the classical warden model of covert communication. A passive warden attempts to decide whether a carrier contains a payload. An active warden applies transformations intended to destroy hidden communication while preserving ordinary utility of the carrier. A forensic adversary may collect many carriers, observe repeated key use, choose messages adaptively, or train a detector directly against a deployed encoder. SRHS is designed primarily for robust covert communication against passive and bounded active wardens; it does not by itself provide payload confidentiality or authenticity, and learned indistinguishability from a trained detector ensemble is not cryptographic security, since empirical detector failure is evidence relative to a declared adversary class, not proof against every possible test.

The same capabilities that support covert communication also support legitimate provenance marking, resilient metadata, synchronization signals, and ownership claims, and they can support malicious coordination or exfiltration. Responsible deployment should include rate controls, explicit authorization for the carriers used, auditable key management, and evaluation by independent steganalysts.

\section{Conclusion}

Steganography has largely been posed as a placement problem: determine where to put a secret while changing the carrier as little as possible. Sparse Recursive Holographic Steganography poses a different problem: determine how to construct a distributed evidence field that remains recoverable across partial observations, and then synthesize that field through a small number of adaptive carrier interventions. Error-correcting projection makes payload evidence global rather than local; sparse routing spends distortion where it carries the highest marginal value; recursive correction lets the encoder observe the realized channel response and revise its allocation without increasing parameter count. None of these components removes the underlying limits of channel capacity, detectability, or adversarial adaptation. If the stated hypotheses hold under the specified ablations, hidden information would no longer be identified with particular altered carrier locations; it would become a recoverable property of the carrier as a whole, weakly present in many views, strengthened by aggregation, and written through adaptive computation rather than fixed placement. The unit of steganographic design would then shift from the secret bit to the secret field, a shift that mirrors, in an adversarial communications setting, the shift Chapter~\ref{ch:ch06-distinction-holonomy} made from the state that persists to the continuation structure that survives transformation: in both cases, the unit worth tracking is not a located point but a distributed, recoverable pattern.
